\section{Potential Impact}

MASFE builds directly on existing NASA land-observation algorithms instead of inventing a new crop model from scratch. The novelty is orchestration: MOD13A1-style EVI screening is cheap enough to run every pass, while MOD11A1-style temperature confirmation is reserved for suspicious patches \cite{mod13a1,mod11a1}. That scheduling choice matters because regenerative agriculture decisions are time-sensitive. Dr.\ Katie Gold's disease-detection work emphasizes that spectral stress often appears before visible symptoms and before whole-field treatment would be agronomically or economically rational \cite{gold-webinar1}.

In simulation, MASFE was tested over 120 passes, 50 field patches, 7 disease events, and 38 cloud-obscured passes. Relative to raw downlink, it reduced transmitted data by 97.8\% and active payload energy by 22.0\%. Relative to a fixed onboard pipeline, it still reduced downlink volume by 4.9\% because it avoids running or transmitting the expensive path when nothing important is happening. The more important metric is mission usefulness: disease-event recall stayed at 100\%, alert precision was 99.0\%, and the false-positive rate remained at 1.0\%.

The farm-scale value proposition is concrete. For a 500 ha regenerative tomato operation, patch-only treatment after a MASFE alert is materially cheaper than blanket preventative spraying and can avoid both unnecessary chemical cost and avoidable yield loss. The same architecture also supports regenerative practices beyond pathogen detection: cover-crop phenology monitoring, soil-carbon MRV support, and no-till verification all benefit from adaptive revisit and selective downlink instead of fixed-rate collection.

The main limitations are honest rather than hidden. The current validation is still TRL 3 and uses synthetic proxy data instead of flight imagery. Cloud screening must be tuned on real scenes, and tropical crop calendars would require policy recalibration. Those gaps are real, but they are the right gaps for a finalist-stage concept: the system question has been answered, and the next work is hardware and data validation rather than rethinking the architecture.
